\documentclass[11pt]{article}
\usepackage[T1]{fontenc}
\usepackage{lmodern,amsmath,amssymb,amsthm}
\usepackage[hmargin=.98in,top=.83in,bottom=.9in]{geometry}
\usepackage[hidelinks]{hyperref}
\linespread{.96}
\newtheorem{theorem}{Theorem}
\newtheorem{lemma}[theorem]{Lemma}
\newtheorem{corollary}[theorem]{Corollary}
\newcommand{\SSer}{\mathfrak S}
\newcommand{\PP}{\mathbb P}
\newcommand{\EE}{\mathbb E}
\newcommand{\one}{\mathbf 1}
\newcommand{\LiK}{\operatorname{Li}}
\title{\vspace*{-.6in}A sifted set with the Hardy--Littlewood tuple counts\\ outside difference two}
\author{\small\href{mailto:admin@cosmosgeek.org}{admin@cosmosgeek.org}
  \qquad\href{mailto:b321014@iiit-bh.ac.in}{b321014@iiit-bh.ac.in}}
\date{Research draft, November 2026}
\begin{document}
\maketitle
\begin{abstract}
We construct a single set of integers containing no pair at distance two.
For every fixed finite tuple of shifts having no such pair, the counting
function has the Hardy--Littlewood main term, also when its starting
point is restricted to any fixed arithmetic progression. In particular
the set has the prime number theorem density in each reduced residue
class. These limiting counts are also stable under deleting both
endpoints of every distance-two pair in a full tuple model.
Our sampling construction gives the stronger finite statement that
all joint probabilities of compatible tuples are preserved exactly.
It uses the fact that the distance-two graph on
integers coprime to six is a matching. Exclusive choices on the two
ends of a matching edge preserve every joint inclusion probability
whose index set does not contain that edge.
\end{abstract}

\section{The counting statement}
How much do the prime-tuple predictions at other patterns say about
the pattern $\{0,2\}$? We allow patterns of every fixed size, also on
each fixed progression of starting points.

Let $H\subset\mathbb Z$ be a finite nonempty set, $k=|H|$, and let
$\nu_p(H)$ count its distinct residue classes modulo $p$. For $q\ge1$
and $a\bmod q$ put
\begin{align}
 \SSer(H;q,a)={1\over q}
 &\prod_{p\mid q}
 {\one_{\{p\nmid a+h\ \text{for all }h\in H\}}\over(1-1/p)^k}
 \prod_{p\nmid q}{1-\nu_p(H)/p\over(1-1/p)^k}.       \label{eq:SS}
\end{align}
For $q=1$ this is the ordinary prime-tuple singular series $\SSer(H)$.
Beyond the primes dividing a nonzero difference in $H$, each local
factor in the second product is $1+O_k(p^{-2})$. Thus the product
converges and is positive whenever every local factor is positive.
Write $\LiK_k(x)=\int_2^x(\log t)^{-k}\,dt$.

\begin{theorem}\label{th:main}
There is a set $A\subset\mathbb N$ with the following properties.
It contains no $u,v$ with $|u-v|=2$. For every $H$ with
$2\notin H-H$, every fixed $q\ge1$ and every $a\bmod q$,
\begin{equation}
 \#\{1\le n\le x:\ n\equiv a\bmod q,\ n+H\subset A\}
 =\SSer(H;q,a)\LiK_k(x)+o_{H,q}\bigl(x/(\log x)^k\bigr).
 \label{eq:theorem}
\end{equation}
If a factor of \eqref{eq:SS} vanishes, only finitely many $n$ in that
progression have $n+H\subset A$.
\end{theorem}

For instance, $\{0,4,10\}$ has its predicted main term, while $H=\{0\}$
gives $\LiK_1(x)/\varphi(q)$ in every reduced residue class.

\paragraph{A deletion comparison.}
Suppose a set $A_0$ satisfies the usual fixed-tuple laws for all
patterns. Delete every element of $A_0$ that has another element at
distance two. For $2\notin H-H$, each lost translate of $H$ requires
a translate in $A_0$ of $H\cup\{h+2\}$ or $H\cup\{h-2\}$ for some
$h\in H$. Each enlarged pattern has $k+1$ elements, so there are
$O_H(x/(\log x)^{k+1})$ losses in total. This argument preserves
the leading counts, also upon restricting the starts to a
progression. The block construction below preserves the full
inclusion probability for each individual compatible configuration,
before any limit or average.

\section{An exclusion that preserves joint probabilities}
Granville's modification of the independent prime model sifts out
small divisors before selecting integers \cite{Granville}. The prime-tuple
statistics of sifted models are discussed by Banks, Ford and Tao
\cite{BFT}. We use fixed forbidden residues and change the sampling
law after sifting. A related probabilistic principle is that joint
occupation probabilities on independent vertex sets determine a
one-dependent hard-core law \cite{HL}. Here every connected component
has at most two vertices and the law is explicit.

\begin{lemma}\label{lem:match}
Let $S$ be any finite subset of the integers coprime to six, and let
$0\le b\le1/2$. There is a random $A_S\subset S$ such that, for every
$U\subset S$,
\[
 \PP(U\subset A_S)=
 \begin{cases}
 b^{|U|}, & 2\notin U-U,\\
 0,       & 2\in U-U.
 \end{cases}
\]
Choices on distinct components of the distance-two graph are independent.
\end{lemma}
\begin{proof}
Both $v-2$ and $v+2$ can be neighbors of $v$ only if $v-2,v,v+2$
all lie in $S$. These three integers meet every residue class modulo
three. Thus each vertex has degree at most one. On each edge, select
its left end alone with probability $b$, its right end alone with
probability $b$, and neither with probability $1-2b$. Select each
isolated vertex with probability $b$. Make all component choices
independently. A set meeting each component at most once has the
claimed product probability. A set containing an edge cannot be selected.
\end{proof}

\section{Local densities on a block}
Use dyadic blocks $[Y,2Y)$, $Y=2^j$, starting at a sufficiently large
$j$. Leave all earlier integers out of $A$. On this block put
\[
 w=\max(3,\log\log Y),\qquad W=\prod_{p\le w}p,\qquad
 b=\frac{W}{\varphi(W)\log Y}.
\]
Retain as candidates the integers coprime to $W$, omitting the first
and last two integers of the block. Apply Lemma~\ref{lem:match} to
the resulting $S$, independently on successive blocks. The omitted
endpoints prevent edges across adjacent blocks. For large $Y$ the
probability is valid, since
$W/\varphi(W)\le\prod_{2\le i\le\lfloor w\rfloor}i/(i-1)
=\lfloor w\rfloor$. Also $W\le w^w=Y^{o(1)}$.

Fix $H,q,a$ as in the theorem. In this section assume $Y$ is large
enough that every prime dividing $q$ is at most $w$. For integer
$0\le t\le Y$ let $N_j(t)$ count starts $Y\le n<Y+t$ with
$n\equiv a\bmod q$, requiring all $n+h$ to be selected within this
block. Allowing patterns to span a block boundary costs
$O_H(1)$ starts per block. Away from its boundaries the probability
is $b^k$ if $(\prod_{h\in H}(n+h),W)=1$, and zero otherwise.

Let $\SSer_w(H;q,a)$ denote \eqref{eq:SS} with the second product
restricted to $p\le w$. The Chinese remainder theorem on a period of
length $L=\operatorname{lcm}(q,W)$ gives the density of candidate
starts before the sampling step as
\[
 {1\over q}\!
 \prod_{p\mid q}\one_{\{p\nmid a+h\ \text{for all }h\in H\}}
 \prod_{\substack{p\mid W\\p\nmid q}}(1-\nu_p(H)/p).
\]
Indeed, modulo $q$ the start has one prescribed class. Each prime of
$W$ outside $q$ removes exactly $\nu_p(H)$ classes independently.
Counting complete periods, bounding the last partial period by $L$,
and multiplying by $b^k$ gives, uniformly in $t$,
\begin{equation}
 \EE N_j(t)=\frac{t}{(\log Y)^k}\SSer_w(H;q,a)+O_{H,q}(W+1).
 \label{eq:expect}
\end{equation}

\section{One set realizing all the counts}
Two translate indicators in $N_j(t)$ are independent whenever they
use disjoint component choices. A nonzero covariance therefore
requires the two starts to differ by an element of
$(H-H)+\{0,2,-2\}$. There are at most $3k^2$ such differences, and
each covariance has absolute value at most one. Consequently
$\operatorname{Var}N_j(t)\ll_k Y$, uniformly in $t$.

For this fixed pattern choose $B=k+2$ and divide $[0,Y]$ into
$\lceil j^B\rceil$ intervals of nearly equal integer length.
By Chebyshev's inequality and the union bound, the probability that
$|N_j(t)-\EE N_j(t)|>Y/j^{k+1}$ at some grid endpoint is
$O_k(j^{B+2k+2}/2^j)$. The series over $j$ converges. By
Borel--Cantelli the bound holds at every grid endpoint for all
sufficiently large $j$, almost surely. Between endpoints, each
additional integer contributes at most one to either the count or
its expectation. The mesh is $O(Y/j^B+1)$, so
\[
 \sup_{0\le t\le Y}|N_j(t)-\EE N_j(t)|
       =o_{H,q}\bigl(Y/(\log Y)^k\bigr)\quad\text{almost surely}.
\]
There are countably many choices of $H,q,a$. We can therefore fix
one realization for which this conclusion holds for every choice.

For each fixed choice, $\SSer_w(H;q,a)\to\SSer(H;q,a)$. On $[Y,2Y]$,
replacing $(\log Y)^{-k}$ by $(\log t)^{-k}$ in the main-term
integral costs $O_k(Y/(\log Y)^{k+1})$ by the mean value theorem.
Equation~\eqref{eq:expect} now gives the desired estimate on full
blocks and arbitrary final partial blocks. Summing the little-oh
errors preserves their order because
$\sum_{j\le J}2^j/j^k=O_k(2^J/J^k)$ after omitting $j=0$.
This proves \eqref{eq:theorem}. If a local factor vanishes at $p$,
all sufficiently large members of $A$ are coprime to $p$, while at
least one entry of every relevant translate is divisible by $p$.
This proves the finite-occurrence assertion, including across block
boundaries.

\paragraph{Scope and provenance.}
The construction concerns sets of integers with the stated counting
laws. Infinitude of ordinary twin primes requires further arithmetic
information about primality. The original Twin Prime project ChatGPT
researcher derived the matching construction, the fixed-shift
marginals and the dyadic variance argument. Codex developed the
simultaneous fixed-tuple and progression statement from those choices
and prepared the exposition. We make no claim of priority.

\begin{thebibliography}{9}
\bibitem{Granville}
A. Granville, \emph{Harald Cram\'er and the distribution of prime
numbers}, Scand. Actuar. J. (1995), suppl.~1, 12--28.
\href{https://doi.org/10.1080/03461238.1995.10413946}{doi:10.1080/03461238.1995.10413946}.
\bibitem{BFT}
W. Banks, K. Ford and T. Tao, \emph{Large prime gaps and probabilistic
models}, Invent. Math. \textbf{233} (2023), 1471--1518.
\href{https://doi.org/10.1007/s00222-023-01199-0}{doi:10.1007/s00222-023-01199-0}.
\bibitem{HL}
A. E. Holroyd and T. M. Liggett, \emph{Finitely dependent coloring},
Forum Math. Pi \textbf{4} (2016), e9.
\href{https://arxiv.org/abs/1403.2448}{arXiv:1403.2448}.
\end{thebibliography}
\end{document}
